\documentclass[12pt,a4paper]{article}

\usepackage[utf8]{inputenc}
\usepackage[T1]{fontenc}
\usepackage[french]{babel}
\usepackage{geometry}
\geometry{top=2.5cm, bottom=2.5cm, left=2.5cm, right=2.5cm}
\usepackage{graphicx}
\usepackage{xcolor}
\usepackage{hyperref}
\usepackage{listings}
\usepackage{fancyhdr}
\usepackage{enumitem}

% Couleurs
\definecolor{primary}{HTML}{1D4ED8}
\definecolor{secondary}{HTML}{0F172A}
\definecolor{lightbg}{HTML}{F8FAFC}
\definecolor{codecolor}{HTML}{0F172A}

\hypersetup{
    colorlinks=true,
    linkcolor=primary,
    urlcolor=primary,
    pdftitle={IncidentFlow - Documentation des Améliorations Non-IA},
}

\pagestyle{fancy}
\fancyhf{}
\lhead{\textbf{IncidentFlow}}
\rhead{Améliorations Fonctionnelles (Non-IA)}
\cfoot{\thepage}

\title{
    \textbf{\Huge IncidentFlow}\\[0.5cm]
    \Large Plateforme Collaborative de Gestion d'Incidents\\[1.5cm]
    \textbf{\Large Rapport d'Implémentation des Améliorations Fonctionnelles et Techniques (Hors-IA)}
    \vspace{1cm}
}
\author{\textbf{Anas HADDOU} \\ Élève Ingénieur en Génie Informatique}
\date{\today}

\begin{document}

\maketitle
\thispagestyle{empty}
\clearpage

\tableofcontents
\clearpage

\section{Introduction}
Dans le cadre de l'optimisation continue de la plateforme \textbf{IncidentFlow}, plusieurs évolutions fonctionnelles et correctifs techniques hors-IA ont été implémentés. L'objectif principal est de renforcer la robustesse de l'application, de fluidifier l'expérience utilisateur (UX) pour les agents de support, et de garantir une gestion rigoureuse de la sécurité et des performances.

Ce document présente les détails d'architecture et d'implémentation de ces quatre améliorations majeures.

\section{Correction et Sécurisation de l'Aperçu PDF}
\subsection{Problématique d'origine}
L'aperçu des pièces jointes PDF au sein de la boîte modale échouait avec une erreur de type \texttt{localhost}. Cet échec était causé par deux facteurs clés :
\begin{itemize}
    \item L'URL de téléchargement était écrite en dur sur \texttt{localhost:8080}, provoquant des blocages d'accès (CORS) ou des erreurs de connexion lorsque l'application était consultée depuis une adresse IP ou un nom de domaine distant.
    \item L'intégration directe d'une URL d'API dans un élément \texttt{<iframe>} ou \texttt{<object>} empêchait la transmission des en-têtes d'authentification requis (token JWT de session et identité de l'utilisateur simulé), renvoyant ainsi une erreur 401 ou 500.
\end{itemize}

\subsection{Solution technique implémentée}
Pour résoudre ces limitations, un flux d'aperçu découplé et authentifié a été mis en œuvre :
\begin{enumerate}
    \item \textbf{Requête Ajax Authentifiée} : Le frontend React effectue un appel \texttt{fetch} en arrière-plan vers l'API d'IncidentFlow en insérant les en-têtes requis :
    \begin{itemize}
        \item \texttt{Authorization: Bearer <token>}
        \item \texttt{X-Mock-User: <email>}
    \end{itemize}
    \item \textbf{Encapsulation par Blob} : La réponse brute de l'API est convertie en objet binaire \texttt{Blob}.
    \item \textbf{Object URL Local} : Une URL d'objet locale temporaire est générée à l'aide de la méthode du navigateur \texttt{URL.createObjectURL(blob)}.
    \item \textbf{Affichage et Nettoyage} : L'URL \texttt{blob:} est passée à l'élément \texttt{<object>} / \texttt{<iframe>}. Le cycle de vie de la modale révoque automatiquement cette URL lors de sa fermeture pour libérer la mémoire vive du navigateur.
\end{enumerate}

\section{Gestion Avancée des Pièces Jointes : Suppression et Renommer}
\subsection{Côté Backend (Spring Boot)}
L'entité \texttt{Attachment} a été dotée de capacités de modification et de cycle de vie complet. Les fonctionnalités suivantes ont été implémentées dans \texttt{IncidentService} :
\begin{itemize}
    \item \textbf{Suppression (\texttt{deleteAttachment})} : L'attachement est supprimé de la base de données relationnelle PostgreSQL, un enregistrement d'action est consigné dans l'historique de l'incident au nom de l'agent effectuant l'opération, et le fichier est physiquement effacé du répertoire \texttt{uploads/} sur le serveur.
    \item \textbf{Renommer (\texttt{renameAttachment})} : Permet la modification logique du nom utilisateur du fichier. Le nom physique reste inchangé (nom unique généré par UUID lors de l'upload) pour garantir l'intégrité du système de fichiers et éviter toute collision.
\end{itemize}
Deux endpoints REST associés ont été configurés dans le contrôleur \texttt{IncidentController} :
\begin{itemize}
    \item \texttt{DELETE /api/incidents/attachments/\{id\}}
    \item \texttt{PUT /api/incidents/attachments/\{id\}/rename}
\end{itemize}

\subsection{Côté Frontend (Vite/React)}
Des boutons d'action d'édition (✏️) et de suppression (🗑️) ont été intégrés à la liste des pièces jointes du panneau latéral. Le renommage est implémenté sous forme d'une zone de saisie textuelle inline (\textit{saisie directe en place}), offrant une transition visuelle de qualité professionnelle.

\section{Éditeur de Commentaires en Markdown Simplifié}
\subsection{Description et Fonctionnalités}
Les techniciens de support ont fréquemment besoin de structurer leurs commentaires pour partager des diagnostics (blocs de code d'erreurs, commandes système, listes d'étapes). Un éditeur Markdown simplifié a été développé, offrant :
\begin{itemize}
    \item Une barre d'outils interactive (Gras, Italique, Liste à puces, Bloc de code de diagnostic).
    \item Un système à double onglet : \textbf{Éditeur} (mode saisie brute) et \textbf{Aperçu} (rendu en temps réel des balises formatées).
\end{itemize}

\subsection{Analyseur de syntaxe Markdown (Parser)}
Un parser Regex ultra-léger et sécurisé a été codé pour convertir le Markdown en HTML structuré. Il intègre un échappement automatique préalable de l'ensemble des balises HTML saisies par l'utilisateur, éliminant tout risque de faille XSS (Cross-Site Scripting).

\section{Compte à Rebours Actif pour les SLAs}
\subsection{Problématique de Performance des Timers}
L'affichage d'un minuteur actif actualisé à la seconde pour chaque incident visible pose des défis de performance majeurs. L'utilisation d'un timer (\texttt{setInterval}) indépendant par incident multiplierait les rafraîchissements React asynchrones et saturerait le CPU et la batterie du client.

\subsection{Architecture Optimisée à Horloge Unique}
Pour éliminer cet inconvénient, l'application utilise une architecture synchronisée à \textbf{horloge globale unique} :
\begin{itemize}
    \item Un unique intervalle racine est démarré au chargement du tableau de bord. Cet intervalle actualise un état central représentant l'horodatage courant exact toutes les secondes.
    \item Les badges d'incidents calculent dynamiquement la différence de temps par rapport à ce signal d'horloge commun lors du rendu général.
    \item \textbf{Impact Design \& UX} : Les couleurs du badge SLA évoluent en temps réel selon la proximité de l'échéance :
    \begin{itemize}
        \item Vert : SLA respecté ou temps restant confortable ($> 30$ min).
        \item Orange : SLA proche (entre 15 et 30 min).
        \item Rouge Pulsant/Clignotant : Alerte critique active (moins de 15 minutes restants) grâce à une animation CSS pulsée fluide, attirant l'attention de l'opérateur.
        \item Rouge sombre : SLA dépassé.
    \end{itemize}
\end{itemize}

\section{Palette de Commandes Universelle (Command Palette)}
\subsection{Description et Raccourci Clavier}
Afin d'offrir une navigation ultra-rapide et l'exécution d'actions directement depuis le clavier, une palette de commandes universelle a été intégrée. Cette barre de recherche globale s'active instantanément à l'aide du raccourci clavier \texttt{Ctrl + K} ou \texttt{Cmd + K} depuis n'importe quel écran de l'application.

\subsection{Implémentation Technique}
\begin{itemize}
    \item \textbf{Écouteur d'Événements Global} : Configuration d'un écouteur global sur l'événement \texttt{keydown} pour intercepter le raccourci clavier et désactiver le comportement par défaut du navigateur (ouverture de la recherche native).
    \item \textbf{Filtrage Intelligent des Commandes et Recherche} :
    \begin{itemize}
        \item Les commandes système sont accessibles via le préfixe \texttt{>} (ex : \texttt{> new} pour déclarer un incident, \texttt{> theme} pour changer de thème, \texttt{> dashboard} pour naviguer).
        \item Sans préfixe, l'utilisateur effectue une recherche floue en direct sur le code de l'incident (ex : \texttt{INC-2026}) ou sur les mots-clés du titre.
    \end{itemize}
    \item \textbf{Navigation Clavier Complète} : L'opérateur peut naviguer entre les résultats à l'aide des flèches directionnelles (\texttt{Up}/\texttt{Down}), valider son choix avec la touche \texttt{Enter} ou fermer la palette via la touche \texttt{Escape}.
\end{itemize}

\section{Journal d'Audit Visuel (Timeline Historique)}
\subsection{Description et Impact Design}
Afin de rendre le suivi des incidents plus lisible et professionnel, le tableau d'historique statique a été transformé en une ligne de temps (timeline) verticale interactive de style Git ou réseau social. Chaque action (changement d'état, assignation, pièce jointe, commentaire) est représentée par une icône et un code couleur spécifiques.

\subsection{Implémentation Technique}
\begin{itemize}
    \item \textbf{Connecteur Vertical Fluide} : Utilisation de bordures CSS absolues pour dessiner le fil conducteur reliant les étapes temporelles de l'incident.
    \item \textbf{Catégorisation et Iconographie Dynamique} : Développement d'un helper React déterminant de manière dynamique la couleur et l'icône de l'événement à partir du texte de l'action :
    \begin{itemize}
        \item \textit{Changement de statut} : Vert avec icône d'activité (\texttt{Activity}).
        \item \textit{Affectation/Assignation} : Violet avec icône d'ajout d'utilisateur (\texttt{UserPlus}).
        \item \textit{Ajout de fichier} : Bleu avec icône d'upload (\texttt{FileUp}).
        \item \textit{Commentaire collaboratif} : Gris ardoise avec icône de message (\texttt{MessageSquare}).
        \item \textit{Création/Déclaration} : Cyan avec icône de plus (\texttt{PlusCircle}).
        \item \textit{Modifications de fichiers (renommage/suppression)} : Jaune ou rouge avec icônes appropriées (\texttt{Edit3}/\texttt{Trash2}).
    \end{itemize}
    \item \textbf{Micro-animations au survol} : Effets de transition doux sur les badges d'icônes et affichage structuré des métadonnées (auteur de l'action et horodatage formaté).
\end{itemize}

\section{Gestion Complète des Commentaires : Modification et Suppression}
\subsection{Description et Règles Métier}
Afin d'offrir une flexibilité maximale aux agents lors des échanges collaboratifs, la possibilité de modifier ou de supprimer un commentaire existant a été ajoutée. Ces actions sont soumises à une règle de sécurité stricte : seul l'auteur d'un commentaire est autorisé à le modifier ou à le supprimer.

\subsection{Implémentation Technique}
\begin{itemize}
    \item \textbf{Côté Backend (Spring Boot)} :
    \begin{itemize}
        \item Injection de \texttt{CommentRepository} pour gérer le cycle de vie direct des commentaires.
        \item Implémentation des méthodes transactionnelles \texttt{updateComment} et \texttt{deleteComment} dans le service, avec vérification systématique de l'identité de l'auteur (\texttt{comment.getAuthor().getId() == currentUser.getId()}).
        \item Exposition de deux endpoints REST sécurisés : \texttt{PUT /api/incidents/comments/\{id\}} et \texttt{DELETE /api/incidents/comments/\{id\}}.
        \item Enregistrement automatique de l'action dans le journal d'historique de l'incident (ex: "Commentaire modifié", "Commentaire supprimé").
    \end{itemize}
    \item \textbf{Côté Frontend (Vite/React)} :
    \begin{itemize}
        \item Ajout des boutons d'actions contextuels (icônes \texttt{Edit3} et \texttt{Trash2}) visibles uniquement pour l'auteur du commentaire (\texttt{currentUser.email == comment.author.email}).
        \item Intégration d'un éditeur inline en place lors de la modification, réutilisant l'éditeur Markdown avec sa barre d'outils et son onglet d'aperçu dynamique.
    \end{itemize}
\end{itemize}

\end{document}
